\section{Countable alphabets by finite projections}
\label{sec:countable-projections}

Let the declared alphabet be \(\mathbb N\). For \(m\geq 0\), let
\(\pi_m\) keep the first \(m\) labels distinct and send all remaining
labels to one tail cell. For a law \(p\) write \(p_m=(\pi_m)_*p\).
Every projection retains its tail cell, including when the reference
assigns that cell probability zero.

\paragraph{Exact fixed-budget equivalence.}
For each nonnegative integer \(n\) and a fixed strict radius \(e\),
a test with \(n\) ordered IID observations exists for \(q\) if and only
if a test with that budget exists for every \(q_m\) at the same radius.
The two separate error bounds are preserved.

To see the forward implication, fix \(m\). Within a cell of positive
\(q_m\)-mass, independently lift a projected observation by drawing
from the conditional law of \(q\) in that cell. In a cell of zero
\(q_m\)-mass use a point mass at an arbitrary declared label of that
cell. These conditionals have disjoint supports. Thus a lifted finite
law \(\widetilde r\) satisfies \(\widetilde q_m=q\) and
\[
 \operatorname{TV}(\widetilde r,q)=\operatorname{TV}(r,q_m).
\]
The full test applied to the independently lifted ordered observations
gives a finite randomized test with the required errors.

For the converse, lift a projected test to a rejection function
\(f_m\colon \mathbb N^n\to[0,1]\) by applying \(\pi_m\) coordinatewise.
The cube \([0,1]^{\mathbb N^n}\), in its product topology, is sequentially
compact. Select a subsequence \(m_j\to\infty\) for which \(f_{m_j}\)
converges pointwise to \(f\colon\mathbb N^n\to[0,1]\). This subsequence
depends only on the sequence of tests. Under each fixed law \(p^{\otimes n}\),
bounded convergence yields
\[
 \mathbb E_{p^{\otimes n}} f
   =\lim_{j\to\infty}\mathbb E_{p^{\otimes n}} f_{m_j}.
\]
The null bound holds at every index. For each \(p\) with
\(\operatorname{TV}(p,q)>e\), the singleton contribution
\(\tfrac12\sum_{i<m}|p_i-q_i|\) converges to
\(\operatorname{TV}(p,q)\). The projected distance is at least this
contribution, so it is strictly greater than \(e\) for all sufficiently
large \(m\). Bounded convergence therefore also gives the stated
alternative bound. This argument includes \(n=0\).

\paragraph{Finiteness at positive radius.}
Pooling the tail \(\mathcal T_m\) loses at most \(q(\mathcal T_m)\) in
TV. In fact, the loss equals
\[
 \frac12\left(\sum_{i\in\mathcal T_m}|p_i-q_i|
                  -|p(\mathcal T_m)-q(\mathcal T_m)|\right)
 \leq \min\{p(\mathcal T_m),q(\mathcal T_m)\}.
\]
For \(e>0\), choose \(m\) with \(q(\mathcal T_m)<e/2\). A finite test for
\(q_m\) at radius \(e/2\) then yields a finite-budget test for \(q\) at
radius \(e\). The use of \(e/2\) proves finiteness only and does not enter
the fixed-budget comparison.

Let \(N_q(e)\) denote the least successful budget at the two fixed error
levels. Fixed-budget equivalence and the option to ignore extra
observations give
\[
 N_q(e)=\sup_{m\geq0}N_{q_m}(e).
\]
In particular, if \(\Phi(q,e)\) is the finite-alphabet formula of the
main theorem, then
\[
 \Phi_{\mathbb N}(q,e):=\sup_{m\geq0}\Phi(q_m,e)
\]
characterizes \(N_q(e)\) with the same universal comparison constants.
This defines a distribution-dependent functional on countable laws,
without a finite runtime guarantee for arbitrary representations of
infinite inputs. Any countably infinite declared alphabet can be
enumerated to use this formulation.

\paragraph{Evaluation directly on a countable null.}
The original head and tail expression also gives a profile on the countable law.
Sort the positive atoms and retain all declared null-zero labels in the tail.
If $e=1$ the alternative is empty. If $e<1$, a countably infinite alphabet
has atoms of arbitrarily small mass and admits a strict alternative.
Only finitely many head atoms are needed to reach $1-e$. Include in the
prefix projection every atom of mass at least the last head mass $a>0$,
and take $m$ large enough that the aggregate unexposed mass $r_m<a$.
All later projections have the same head probabilities, $P$, $w$ and
$\Delta$. This also covers finite positive support on an infinite alphabet,
for which a zero-mass tail cell remains available once that support is exposed.

For fixed $n,x>0$, the contribution of a single atom $z$ to $K(n,x)$
is superadditive in $z$ since its ratio to $z$ is nondecreasing.
The same holds for the summands of $J$. Consequently $K_m(n,x)$ and
$J_m$ decrease to the corresponding formulae on the full sequence, with
errors at most $r_m/x$ and $r_m$ respectively. If $w>0$, the continuous
strictly decreasing function $K(n,\cdot)$ has a unique positive unit level.
For every $t>X(n)$, eventually $K_m(n,t)<1$ whereas $K_m(n,X(n))\ge1$.
Thus $X_m(n)\downarrow X(n)$. When $w=0$ all roots vanish. The small-overlap
branch stabilizes even if $J=16\delta$, as $J_m\ge J$. Its rounded reciprocal
stabilizes from below. On the regular branch the first crossing $n_*$ for
the full sequence is finite by the same full-head bound and $Y(1)>0$.
It passes every sufficiently fine projected comparison even at equality.
Every earlier dyadic comparison fails strictly in the limit, so all the
finitely many earlier comparisons eventually fail. Hence $\Phi(q_m,e)$
eventually equals the direct countable formula, while $N_{q_m}(e)$ is a
bounded nondecreasing integer sequence and eventually equals $N_q(e)$.
The finite theorem therefore gives the same comparison constants for the
direct profile. At any positive tail root the saturated prefix in
\eqref{eq:radical} is finite. A countable infimum describes the minimal
tail expression at fixed $n$. No finite representation is assumed for
an arbitrary infinite probability law.
